
\documentclass[aspectratio=169]{beamer}

\mode<presentation>
\usetheme{Boadilla}
\definecolor{NTHU1}{RGB}{127,16,132}
\definecolor{NTHU2}{RGB}{228,210,231}
\definecolor{blue}{RGB}{30,90,205}
\definecolor{red}{RGB}{213,94,0}
\definecolor{green}{RGB}{0,128,0}
\definecolor{charcoalgray}{RGB}{108,104,104}
\setbeamercolor{title}{fg=NTHU1}
\setbeamercolor{frametitle}{fg=NTHU1}
\setbeamercolor{block title}{bg=NTHU1, fg=white}
\setbeamercolor{block body}{bg=NTHU2}
\setbeamercolor{structure}{fg=NTHU1}
\setbeamercolor{item projected}{fg=white}
\setbeamercolor{item}{fg=NTHU1}
\setbeamercolor{subitem}{fg=NTHU1}
\setbeamercolor{section in toc}{fg=NTHU1}
\setbeamercolor{description item}{fg=NTHU1}
\setbeamercolor{caption name}{fg=NTHU1}
\setbeamercolor{button}{bg=NTHU1, fg=white}
\setbeamercolor{caption name}{fg=NTHU1}
\usepackage{graphics}
\usepackage{tikz}
\usepackage{amsmath}
\usepackage{bbm}
\usetikzlibrary{decorations.pathreplacing}
\usepackage{geometry}
\usepackage{booktabs}
\usepackage{bookmark}
\usepackage{multirow, makecell}
\usepackage{float}
\usepackage{fancyvrb}
\usepackage{caption}
\captionsetup[figure]{singlelinecheck = false, format= hang, justification=centering}
\captionsetup[subfigure]{singlelinecheck = false, format= hang, justification=raggedright}
\usepackage{subcaption}
\usepackage{adjustbox}
\usepackage{hyperref}
\usepackage{threeparttable}
\usepackage{appendixnumberbeamer} 
\usepackage[T1]{fontenc}
\usepackage[default]{lato} 
\usepackage{smartdiagram}
\smartdiagramset{
  back arrow disabled = true,
  module minimum width=1cm,
  module x sep = 3,
  uniform arrow color=true,
}
\usepackage{xcolor}
\usepackage{colortbl}
	\definecolor{light-gray}{gray}{0.99}

\newenvironment{wideitemize}{\itemize\addtolength{\itemsep}{10pt}}{\enditemize}
\newenvironment{wideenumerate}{\enumerate\addtolength{\itemsep}{10pt}}{\endenumerate}
\newenvironment{widedescription}{\description\addtolength{\itemsep}{10pt}}{\enddescription}
\hypersetup{
colorlinks=true,
linkcolor=NTHU1,
filecolor=green, 
urlcolor=blue,
}
\beamertemplatenavigationsymbolsempty
\setbeamercolor{author in head/foot}{bg=white, fg=NTHU1}
\setbeamercolor{title in head/foot}{bg=white, fg=NTHU1}
\setbeamercolor{date in head/foot}{bg=white, fg=NTHU1}
\setbeamercolor{section in head/foot}{bg=white, fg=NTHU1}
\setbeamercolor{headline}{bg=white}
\setbeamertemplate{footline}{
    \leavevmode%
    \hbox{%
        \begin{beamercolorbox}[wd=.333333\paperwidth,ht=2.25ex,dp=1ex,center]{date in head/foot}%
            \usebeamerfont{date in head/foot}%\insertdate
        \end{beamercolorbox}%
        \begin{beamercolorbox}[wd=.333333\paperwidth,ht=2.25ex,dp=1ex,center]{title in head/foot}%
            \usebeamerfont{title in head/foot}\insertshorttitle
        \end{beamercolorbox}%
        \begin{beamercolorbox}[wd=.333333\paperwidth,ht=2.25ex,dp=1ex,right]{date in head/foot}%
            \usebeamerfont{date in head/foot}\insertframenumber{}\hspace*{2em}
        \end{beamercolorbox}}%
        \vskip0pt%
    }
    %% May need to script the introduction!!!!!!
%\setbeamercolor{page number in head/foot}{fg=NTHU1}
\setbeamertemplate{section in toc}[sections numbered]
\setbeamertemplate{subsection in toc}{\leavevmode\leftskip=3em\rlap{\hskip-1.75em\inserttocsectionnumber.\inserttocsubsectionnumber}
\inserttocsubsection\par}
\setbeamerfont{subsection in toc}{size=\footnotesize}
%\setbeamertemplate{headline}{%
  %\begin{beamercolorbox}[ht=5.5ex]{section in head/foot}
    %\vskip2pt\insertnavigation{0.33\paperwidth}\vskip2pt
  %\end{beamercolorbox}%
%}
\newenvironment{transitionframe}{\setbeamercolor{background canvas}{bg=white}\setbeamertemplate{footline}{} \begin{frame}[noframenumbering]}{\end{frame}}


\makeatletter
\let\@@magyar@captionfix\relax
\makeatother


\title[]{Public Economics and Development} % Change this regularly
\subtitle[]{Week 3: Inequality across time and countries}
\author[Lee]{Seung-hun Lee }
\institute[]{Taipei School of Economics\\ National Tsing Hua University}

\date[]{March 10th, 2026}

\begin{document}
\begin{transitionframe}
\titlepage
\end{transitionframe}


%%% Color slides for section headers: Use for colloquium version (The ones bewteen \iffals and \fi)

\section{Overview of the topic}
\begin{frame}
\frametitle{What is inequality?} 
Unequal distribution of income, wealth, and opportunities
\begin{itemize}\setlength\itemsep{3pt}
\item Income: Differences in flow of money (e.g. labor earnings)
\item Wealth: Differences in value (stock) of assets accumulated over time (stock, homes etc)
\item Opportunities: Life chances determined by circumstances beyond individual effort
\end{itemize}\medskip
How are they measured?
\begin{itemize}\setlength\itemsep{3pt}
\item Classical: Gini Coefficient, Lorenz curve, ratio of top 20\% over bottom 80\% etc
\item Problem: Large inequality even within the top
\begin{itemize}\setlength\itemsep{3pt}
\item Increasing share of wealth and income  by top 1\%, skewing classical measurements further

\end{itemize}
\item Increasing focus on share of income and wealth of top 10\%, 1\% and 0.1\%  
\begin{itemize}\setlength\itemsep{3pt}
\item Very few people take up significant share of income and taxes
\end{itemize}
\item Problem: Income measures may not be consistent 
\begin{itemize}\setlength\itemsep{3pt}
\item Different tax laws, income-earning opportunities, social contexts, and data availability
\end{itemize}
\end{itemize}
\end{frame}

\begin{frame}
\frametitle{We are even talking about top 0.01\% nowadays}
\begin{figure}[H]
  \centering
  \includegraphics[keepaspectratio, width=1\textwidth]{fig1.png}
\end{figure}
\begin{small}
Source: ``Never mind the 1 percent. Let's talk about 0.01 percent'', \textit{Chicago Booth Review} \href{https://www.chicagobooth.edu/review/never-mind-1-percent-lets-talk-about-001-percent}{(link)}
\end{small}
\end{frame}

\begin{frame}
\frametitle{Why does it matter?} 


     Efficiency losses: Inequality constrains investment in human capital
    \begin{itemize}
        \item Talented but poor individuals cannot invest in education or businesses
        \item Aggregate productivity suffers even if top earners are highly productive
    \end{itemize}\medskip
     Political economy: Wealthy groups gain outsized influence over policy
    \begin{itemize}
        \item Redistribution becomes harder precisely when it is most needed
        \item May even lead to lack of accountability on those with power
        \item Risk of self-reinforcing cycle: inequality $\to$ political power $\to$ more inequality
    \end{itemize}\medskip
     Social cohesion: High inequality linked to lower trust, worse health, higher crime
    \begin{itemize}
        \item Effects felt across the income distribution, not just at the bottom
    \end{itemize}\bigskip
    Ultimately, Rising inequality $\Rightarrow$ concentration of opportunities and resources on few people
\end{frame}


\begin{frame}
\frametitle{How are inequality passed on?}
Across generations 
\begin{itemize}\setlength\itemsep{3pt}
    \item Parental income and wealth strongly predict children's economic outcomes
    \begin{itemize}\setlength\itemsep{3pt}
      \item Direct wealth transfers: bequests and inheritances  
      \item Better schooling, networks, and risk-taking capacity 
    \end{itemize}
    \item Thus, many cases where wealth begets wealth
\end{itemize} \medskip
Across different groups within same generation
\begin{itemize}\setlength\itemsep{3pt}
 \item Technology reshapes labor demand, favoring some workers over others
    \begin{itemize}
        \item Automation displaces workers with skills substitutable with technology 
    \end{itemize}
    \item Returns to education and skills have risen, widening the college wage premium
    \item Dispersed resources also explains inequality across gender, regions, race and etc
\end{itemize}
\end{frame}

\begin{frame}
\frametitle{Income distribution across education attainment (US example)}
\begin{figure}[H]
  \centering
  \includegraphics[keepaspectratio, width=0.9\textwidth]{fig2.png}
\end{figure}
\end{frame}

\begin{transitionframe}
\frametitle{Roadmap for today}
\tableofcontents
\end{transitionframe}


\section{Measuring Inequality}
\subsection{Atkinson, Piketty, Saez (2011): The Top Incomes in the Long Run of History}
\begin{transitionframe}
\begin{center}
  \begin{huge}
    \textcolor{NTHU1}{%
  Atkinson, Piketty, Saez (2011): \\ The Top Incomes in the Long Run of History
    }
  \end{huge}
\end{center}
\end{transitionframe}


\begin{frame}
\frametitle{Q: What can we learn from top income shares} 
Do top income shares matter?
\begin{itemize}\setlength\itemsep{3pt}
\item If highly unequal, command of resource and growth concentrated at the top
\item Sizeable impact on \textit{overall} inequality 
\item Thus, top income shares carry a sizable implications for policy designs
\end{itemize}\medskip
This paper surveys trends in top incomes and discusses the methodology
\begin{itemize}\setlength\itemsep{3pt}
\item Identifies differential trends in top shares across the globe
\item Explores what factors drive the trends in top income shares
\item Faces challenges regarding consistency of income definition across time and globe
\end{itemize}
\end{frame}

\begin{frame}
\frametitle{Top income shares generally show a U-shape}
\begin{figure}[H]
  \centering
  \includegraphics[keepaspectratio, width=0.75\textwidth]{aps(2011)_f1.png}
\end{figure}
\end{frame}

\begin{frame}
\frametitle{With different trends in some places}
\begin{figure}[H]
  \centering
  \includegraphics[keepaspectratio, width=0.75\textwidth]{aps(2011)_f2.png}
\end{figure}
\end{frame}

\begin{frame}
\frametitle{Primarily driven by labor earnings and capital gains}
\begin{figure}[H]
  \centering
  \includegraphics[keepaspectratio, width=0.8\textwidth]{aps(2011)_f3.png}
\end{figure}
\end{frame}

\begin{frame}
\frametitle{Challenges faced in identifying top income shares} 
Challenges faced within the same country
\begin{itemize}\setlength\itemsep{3pt}
\item How many tax filers? Family + single individuals $\to$ Difficult to trace individuals
\item What is income? Definition differs across time and jurisdictions
\begin{itemize}
\item Especially true for capital gains (profits from sale of capital goods and assets)
\item Some countries treat it as income, others do not
\end{itemize}
\item Past data: Not available at all or with significant gaps
\end{itemize}\medskip
Other potential problems
\begin{itemize}\setlength\itemsep{3pt}
\item Top income possibly underreported
\item This is due to tax avoidance (legal) or tax evasion (illegal)
\end{itemize}
\end{frame}


\begin{frame}
\frametitle{Where are we now in our understanding of top shares?} 
Widespread usage of tax data then before
\begin{itemize}\setlength\itemsep{3pt}
\item One source with long time periods and wide coverage of population
\item Useful for distributional analysis (despite limits on point-estimates)
\item More economics papers using tax data to study income and wealth
\end{itemize}\medskip
What needs to be done
\begin{itemize}\setlength\itemsep{3pt}
\item Inclusion of developing countries (many still missing)
\item Complement tax data with inherited wealth or company data
\item Consistent defintion of income that takes into account well-defined tax units and income categories
\end{itemize}\medskip
\end{frame}

\subsection{Saez, Zucman (2016): Wealth Inequality in the United States Since 1913: Evidence from Capitalized Income Tax Data}
\begin{transitionframe}
\begin{center}
  \begin{huge}
    \textcolor{NTHU1}{%
  Saez, Zucman (2016): \\ Wealth Inequality in the United States Since 1913: Evidence from Capitalized Income Tax Data
    }
  \end{huge}
\end{center}
\end{transitionframe}


\begin{frame}
\frametitle{Q: How did inequality of wealth in US change over time?}
Currently available evidence on wealth inequality is mixed
\begin{itemize}\setlength\itemsep{3pt}
\item Survey of Consumer Finances, SCF: top 1\% has 35\% of wealth in 2013, (30\% in 1992)
\item Estate tax returns: top 1\% has 20\% of wealth in early 2010s (less than DEN, FIN, FRA)
\item While wealth inequality has risen since 1970s, no concensus on exact figures
\end{itemize}\medskip
This paper sheds new light on evolution of wealth inequality with income tax data from IRS
\begin{itemize}\setlength\itemsep{3pt}
\item Based on capital income reported by taxpayers (dividens, interest, rents, profits etc)
\item Computes capitalization factor mapping flow of tax income to amount of wealth recorded in household balance sheets
\item Creates long-run series of wealth distribution consistent with macroeconomic totals
\item Finds that top 1\% owns 42\% of wealth in 2012, driven by rise of top 0.1\% shares
\end{itemize}\medskip
\end{frame}

\begin{frame}
\frametitle{What is wealth?}
Definition of wealth used in the research
\begin{itemize}\setlength\itemsep{3pt}
\item Current market value of all household-owned assets net of debts
\item Most pension wealth, and all (non-)financial assets where ownership is clearly defined
\begin{itemize}
\item Not included: Human capital, nonprofit institutions wealth (no clear market value), durables (no tax return info)
\end{itemize}
\item Aggregate wealth of US: Based on household balance sheets of US Financial Accounts
\begin{itemize}
\item Obtained from SCF, IRS, US Treasury, BEA, Census, Federal Housing Finance Agency etc
\end{itemize}
\end{itemize}\medskip
From income to wealth
\begin{itemize}\setlength\itemsep{3pt}
\item Start with capital income (9 categories) reported on individual tax returns 
\item Then, capitalize these incomes, while excluding components without taxable income
\item Rise in capital inequality: Top 0.1\% share went from 10\% in 1960s  $\to$ 33\% in 2012
\end{itemize}\medskip
\end{frame}

\begin{frame}
\frametitle{Capitalizing the income}
How capitalization technique works
\begin{itemize}\setlength\itemsep{3pt}
\item Map each capital income category to wealth category in household balance sheet
\item For each category, compute annual capitalization factor: $\frac{\text{Aggregate HH wealth}}{\text{Tax return income}} $ 
\begin{itemize}\setlength\itemsep{3pt}
\item Understand it as how much wealth is required to generate a dollar of income
\item Capitalization factor $\uparrow\to$ Assets that generate relatively small income (savings account)
\end{itemize}
\item Then, exclude assets that does not generate taxable income
\begin{itemize}\setlength\itemsep{3pt}
\item Owner occupied housing (no rents), life insurance, offshore wealth etc
\end{itemize}
\end{itemize}\medskip
Potential challenges
\begin{itemize}\setlength\itemsep{3pt}
\item Assets generate different rate of income across households and level of wealth
\item Fortunately, this method matches well with wealth held by foundations and SCF data
\end{itemize}\medskip
\end{frame}

\begin{frame}
\frametitle{Wealth share of Top 0.1\% = That of bottom 90\%}
\begin{figure}[H]
  \centering
  \includegraphics[keepaspectratio, width=0.75\textwidth]{sz(2016)_t1.png}
\end{figure}
\end{frame}



\begin{frame}
\frametitle{Rising wealth inequality in the past 40 years}
\begin{figure}[H]
  \centering
  \includegraphics[keepaspectratio, width=0.8\textwidth]{sz(2016)_f1.png}
\end{figure}
\end{frame}

\begin{frame}
\frametitle{Driven by surge in the wealth share of top 1\%}
\begin{figure}[H]
  \centering
  \includegraphics[keepaspectratio, width=0.8\textwidth]{sz(2016)_f2.png}
\end{figure}
\end{frame}

\begin{frame}
\frametitle{Driven by rising share of inequality in saving rates}
\begin{figure}[H]
  \centering
  \includegraphics[keepaspectratio, width=0.8\textwidth]{sz(2016)_f3.png}
\end{figure}
\end{frame}
 
\subsection{Auten, Splinter (2024): Income Inequality in the United States: Using Tax Data to Measure Long-Term Trends}
\begin{transitionframe}
\begin{center}
  \begin{huge}
    \textcolor{NTHU1}{%
  Auten, Splinter (2024):\\ Income Inequality in the United States: Using Tax Data to Measure Long-Term Trends
    }
  \end{huge}
\end{center}
\end{transitionframe}


\begin{frame}
\frametitle{Q: How unequal are income when measured with consistent data} 
Measure of income distribution over time is difficult
\begin{itemize}\setlength\itemsep{3pt}
\item Different social conditions across time (marriages, household size etc)
\item Different education: Join labor market late but earn more
\item Business cycles and inflation rates differ across periods
\item Tax rules also change across time, and share of income missing in tax data have risen
\end{itemize}\medskip
This paper presents new estimates of the trends in income distribution (from tax data)
\begin{itemize}\setlength\itemsep{3pt}
\item With accurate measures of income taking tax reforms and technicalities into account
\item Social changes are also taken into account (e.g. household size and marriage rates)
\item Lower top income shares and less increase in inequality compared to other works
\end{itemize}\medskip
\end{frame}

\begin{frame}
\frametitle{How are tax data used here?}
 Start with a seminal inequality measure from Piketty and Saez
\begin{itemize}\setlength\itemsep{3pt}
\item Create an \textit{improved} fiscal income adjusting for tax changes in 1986
\item Also take family structures (marriage rates) and sample issues into account
\item Pre-tax income: Follows national income concept ()
\item Pre-tax + transfers: includes government transfers, more complete economic measure
\item After-tax: deducts federal, state, local taxes
\end{itemize}\medskip
Why tax reform in 1986 and other changes matter
\begin{itemize}\setlength\itemsep{3pt}
\item Top income tax rate $<$ corporate tax rate 
\item So some business income shows up as individual incomes
\item Broadening of tax base (those subject to taxation): Dependent and nonresident filers $\uparrow$
\item Poor less likely to marry $\to$ total tax unit at bottom $\uparrow\to$ Income understated at bottom
\end{itemize} \medskip
The authors adjust these factors and re-estimate income distribution across time
\end{frame}

\begin{frame}
\frametitle{Income distribution of bottom 50\% more optimistic}
\begin{figure}[H]
  \centering
  \includegraphics[keepaspectratio, width=1\textwidth]{as(2024)_f1.png}
\end{figure}
\end{frame}


\begin{frame}
\frametitle{Different top 1\% estimates}
\begin{figure}[H]
  \centering
  \includegraphics[keepaspectratio, width=1\textwidth]{as(2024)_f2.png}
\end{figure}
\end{frame}

\begin{frame}
\frametitle{Different top 0.1\% estimates}
\begin{figure}[H]
  \centering
  \includegraphics[keepaspectratio, width=0.6\textwidth]{as(2024)_f3.png}
\end{figure}
\end{frame}

\begin{frame}
\frametitle{Where do the differences occur?}
1/3 of increase in top income occured around major tax reforms 
\begin{itemize}\setlength\itemsep{3pt}
\item Lower tax rates and broader tax base $\to$ Changing incentives to report income
\item Top income shares in 1960s-70s rise after adjustment (i.e incomes were `sheltered')
\end{itemize} \medskip
Note that married couples file together
\begin{itemize}\setlength\itemsep{3pt}
\item Lower marriage rates $\to$ more tax units $\to$ more high income tax units at top
\item This leads to overestimated top incomes in recent years
\end{itemize}\medskip
Other issues
\begin{itemize}
 \item Different treatment of pensions and business income that accrue over multiple years
 \item Drop non-resident and dependent filers
\end{itemize}
\end{frame}


\begin{frame}
\frametitle{Takeaways and remaining questions} 
Income distribution trends
\begin{itemize}\setlength\itemsep{3pt}
\item Top 1\% income share growth are modest compared to other estimates
\item Bottom 50\% income share grew more compared to previous studies
\item Current estimates incorporate transfers, tax reforms, and social changes across time
\end{itemize}\medskip
So what to make of these differences?
\begin{itemize}\setlength\itemsep{3pt}
\item Data discrepencies and different assumption about missing income could lead to significant differences in estimates
\item This is general in any empirical studies: Lots of judgement calls involved
\item Conclusions are very sensitive to these choices
\item These choices should be defensible from reasonable researchers' perspective
\end{itemize}\medskip
\end{frame}








\section{Transmission of inequality}
\subsection{Black, Devereux, Lundborg, Majlesi (2020): Poor Little Rich Kids? The Role of Nature versus Nurther in Wealth and Other Economic Outcomes and Behaviors}
\begin{transitionframe}
\begin{center}
  \begin{huge}
    \textcolor{NTHU1}{%
  Black, Devereux, Lundborg, Majlesi (2020):\\ Poor Little Rich Kids? The Role of Nature versus Nurther in Wealth and Other Economic Outcomes and Behaviors
    }
  \end{huge}
\end{center}
\end{transitionframe}

\begin{frame}
\frametitle{Q: What determines correlation of wealth across generation?} 
Is it biology (nature) or environment (nurture)?
\begin{itemize}\setlength\itemsep{3pt}
\item Nature: Genetic inheritance of skills, attitudes, and preference
\begin{itemize}\setlength\itemsep{3pt}
\item Intergenerational correlations arise because wealthy families are inherently more talented
\item They would be wealthier than others even without growing up with wealthier parents 
\end{itemize}
\item Nurture: Parental investments in education, opportunities, and preference \begin{itemize}\setlength\itemsep{3pt}
\item Intergenerational correlations arise by environment the children grows up in
\item Any children given these opportunities may benefit
\end{itemize}
\item These two can interact: Environmental effect depending on biological endowments
\end{itemize} \medskip
This study disentagles nature vs nurture in intergenerational wealth using adoptees
\begin{itemize}\setlength\itemsep{3pt}
\item Observe how wealth of adoptive children relate to biological and adoptive parents
\item Use Swedish adminstrative data covering adopted children born in 1950-70
\item Are wealth diffrent from education, income, risk preferences? % in terms of intergenerational transmission
\end{itemize}
\end{frame}

\begin{frame}
\frametitle{Data sources}
 Swedish administrative data with all citizens born between 1932-80
\begin{itemize}\setlength\itemsep{3pt}
\item Education attainment, residence, other demographics
\item Swedish-born adoptees with biological and adoptive parents
\item Wealth register data (used for taxes): Various financial assets, debt, real assets
\item Also: Income, labor earnings, savings, and consumption (clocked at 2006)
\end{itemize}\medskip
 Selection of samples
\begin{itemize}\setlength\itemsep{3pt}
\item Children born in 1950-70: Avoid someone too young to accumulate wealth
\item At least one parent alive in 2006: Focus on pre-bequest wealth
\item 1.2 million children, and 2,598 adopted children
\item Problem: Adoptive parents relatively richer than general ($\because$ adoption requirements) 
\end{itemize}\medskip
\end{frame}

\begin{frame}
\frametitle{Empirical Strategy on intergenerational transmission}
Correlate within-cohort ranks of net wealth of adoptees and parents
\[
W_{ij}=\beta_0 + \beta_1 W_i + \beta_2 W_j + \beta_3 W_i^2 + \beta_4 W_j^2 + \beta_5 W_i W_j + X\beta_6 + e_{ij}
\]\vspace{-15pt}
\begin{itemize}\setlength\itemsep{3pt}
\item $W$ stands for within-cohort ranks of wealth (left: children, right: parents)
\item $i$ indexes biological and $j$ indexes adoptive parents
\item $X$ are covariates: Gender, region of residence, cohort dummies for parent and child
\item Capture linearity ($\beta_1,\beta_2$), nonlinearity ($\beta_3,\beta_4$) and interactions ()$\beta_5$) 
\end{itemize}\medskip
Identifying assumption
\begin{itemize}\setlength\itemsep{3pt}
\item Random assignment of adoptive families at birth
\item i.e. adoption not based on systematic matching on family and child characteristics
\end{itemize}\medskip
\end{frame}

\begin{frame}
\frametitle{Wealth predominantly associated with adoptive parents}
\begin{figure}[H]
  \centering
  \includegraphics[keepaspectratio, width=0.7\textwidth]{bdlm(2020)_t1.png}
\end{figure}
\end{frame}


\begin{frame}
\frametitle{And largely in a linear relationship}
\begin{figure}[H]
  \centering
  \includegraphics[keepaspectratio, width=0.7\textwidth]{bdlm(2020)_f1.png}
\end{figure}
\end{frame}

\begin{frame}
\frametitle{Environment matters for wealth, Nature matters for some others}
\begin{figure}[H]
  \centering
  \includegraphics[keepaspectratio, width=0.7\textwidth]{bdlm(2020)_t2.png}
\end{figure}
\end{frame}

\begin{frame}
\frametitle{Takeaways and remaining questions} 
Key findings
\begin{itemize}\setlength\itemsep{3pt}
\item Wealth transmission is primarily due to environment
\item Not because children of wealtheir famlies are more inherently able
\item Environment determines other wealth-related variables (biology for human capital)
\item No meaningful interaction between nature and nurture
\end{itemize}\medskip
Remaining questions
\begin{itemize}\setlength\itemsep{3pt}
\item Wealth begets wealth: So how does this relate to growing inequality?
\item Poorer familes increasingly have fewer opportunities and environmental advantages 
\item Need for policy to equalize opportunities and mitigate intergenerational disparities
\end{itemize}
\end{frame}



\subsection{Autor, Levy, Murnane (2003): The Skill Content of Recent Technological Change: An Empirical Exploration}
\begin{transitionframe}
\begin{center}
  \begin{huge}
    \textcolor{NTHU1}{%
Autor, Levy, Murnane (2003):\\ The Skill Content of Recent Technological Change: An Empirical Exploration
    }
  \end{huge}
\end{center}
\end{transitionframe}


\begin{frame}
\frametitle{Q: How does computerization affect job skill demands?} 
Why are educated workers more in demand than lesser-educated workers?
\begin{itemize}\setlength\itemsep{3pt}
\item Striking correlation between adoption of computers and demand for college graduates
\item Rise in earnings gap between college and non-college graduates (skill premium)
\item But what about computers (or what people do with computers) explains this pattern?
\end{itemize}\medskip
This paper: Computers change tasks workers perform, affecting demand for skills 
\begin{itemize}\setlength\itemsep{3pt}
\item Computers complement certain human tasks, while substituting other human tasks
  \begin{itemize}\setlength\itemsep{3pt}
\item Routine tasks: Cognitive and manual tasks that follow explicit guidelines
\item Non-routine tasks: Problem-solving and complex communication activities
\end{itemize}
\item Computers substitute for routine tasks, complement non-routine cognitive tasks
\item Notes shfits away from routine tasks to non-routine ones, regardless of education
\end{itemize}\medskip
\end{frame}

\begin{frame}
\frametitle{What types of tasks?}
Tasks are categorized based on what machines can do (and can't) 
\begin{itemize}\setlength\itemsep{3pt}
\item Can perform repetitive tasks with instructions accurately and rapidly
\item They can also perform symbolic processing (storing, retriving information)
\end{itemize}\medskip
Classification of tasks
\begin{itemize}\setlength\itemsep{3pt}
\item Routine tasks can be accomplished with machines, based on programmed rules
\begin{itemize}\setlength\itemsep{3pt}
\item Many manual tasks (not all) fall into this category
\item But some cognitive ones too: Calculating, coordinating, record-keeping
\end{itemize}
\item Non-routine tasks require flexibility, creativity, problem-solving, and communications
\begin{itemize}\setlength\itemsep{3pt}
\item Often occurs without explicit guidelines
\end{itemize}
\item Prediction: Computers replace routine tasks but complement non-routine works
\end{itemize}
\end{frame}

\begin{frame}
\frametitle{Some examples} 
\begin{figure}
  \centering
  \includegraphics[keepaspectratio, width=0.67\textwidth]{alm(2003)_t1.png}
\end{figure}
\end{frame}


\begin{frame}
\frametitle{Confirming this with data}
Data source
\begin{itemize}\setlength\itemsep{3pt}
\item Dictionary of Occupational Titles (DOT) from US Department of Labor
\item Evaluate intensity of routines for 12,000 occupations across 44 dimensions
\item Merged with Census (based on occupation codes) for individual level characteristics 
\end{itemize}\medskip
Empirical strategy: Measuring changing job task requirements over time
\begin{itemize}\setlength\itemsep{3pt}
\item Changes over time in occupational distribution of employment (extensive)
\item Also keep track of task contents within occupation (intensive)
\item Also tracks differences in computer adoptaion across industries and time
\end{itemize}
\end{frame}

\begin{frame}
\frametitle{Trends in occupational distribution based on task content} 
\begin{figure}
  \centering
  \includegraphics[keepaspectratio, width=0.64\textwidth]{alm(2003)_f1.png}
\end{figure}
\end{frame}


\begin{frame}
\frametitle{Changes in tasks within industries} 
\begin{figure}
  \centering
  \includegraphics[keepaspectratio, width=0.575\textwidth]{alm(2003)_t2.png}
\end{figure}
\end{frame}


\begin{frame}
\frametitle{Takeaways and remaining questions}
Takeaways
\begin{itemize}\setlength\itemsep{3pt}
\item Computerization substitutes for routine tasks and complements non-routine ones
\item This explains a decrease in demand for routine tasks since 1970s
\item Tasks, not just skills accumulated through education, explains gap in demand
\end{itemize}\medskip
Remaining questions: Note that this paper came out in 2003
\begin{itemize}\setlength\itemsep{3pt}
\item AI introduces a new challenge: Evidence that non-routine tasks are substitutable
\item Thus, it is feasible that demand for these workers can change
\item Updated explanations for earnings gap and inequality in demand required

\end{itemize}
\end{frame}

\subsection{Acemoglu, Resrepo (2022): Tasks, Automation, and the Rise in US Wage Inequality}
\begin{transitionframe}
\begin{center}
  \begin{huge}
    \textcolor{NTHU1}{%
Acemoglu, Resrepo (2022):\\ Tasks, Automation, and the Rise in US Wage Inequality
    }
  \end{huge}
\end{center}
\end{transitionframe}


\begin{frame}
\frametitle{Q: How does automation affect wage inequality?} 
Rising wage inequality across industrialized economics 
\begin{itemize}\setlength\itemsep{3pt}
\item Real wages of post-college graduates on the rise
\item Those of workers with lower educational attainment stagnant or falling
\item Could this be explained by differential displacement following automation?
\begin{itemize}
\item Workers with routine tasks vs tasks not replaceable with automation
\end{itemize}
\end{itemize}\medskip
This paper provides and tests the framework explaining how automation affects wages
\begin{itemize}\setlength\itemsep{3pt}
\item Lays out channels for how technological changes affect wages
\item It can displace workers, and also reduce real wages of workers not yet displaced
\item Empirically investigate relationship between wage changes and task displacement across different groups
\end{itemize}
\end{frame}

\begin{frame}
\frametitle{Link between automation and displacement: Very simplified} 
Allocation of task $x$ across capital and different types of labor
\begin{itemize}\setlength\itemsep{3pt}
\item Each task (say knitting clothes) is allocated across capital and different types of labor
\item There is different productivity specific to each factor of production
\item And the productivity of a given factor can differ based on tasks it performs
\end{itemize}\medskip
Technology can affect productivity and factors in these ways
\begin{itemize}\setlength\itemsep{3pt}
\item It can augment productivity of capital / labor as a whole
\item It can increase productivity of a factor doinbg a particular task
\item It can increase productivity of capital at tasks previously carried out by labor
\begin{itemize}
\item The last case leads to displacement of labor: Robots replacing blue-collar tasks
\end{itemize}
\item Technologies improving productivity of skilled labor do not involve displacement
\item Technologies that reduce comparative advantage of workers leds to displacement
\end{itemize}\medskip
\end{frame}


\begin{frame}
\frametitle{Data and empirical strategy} 
Data sources
\begin{itemize}\setlength\itemsep{3pt}
\item BEA: industry labor shares, factor prices, and value added
\item Census and American Community Survey: labor market outcomes by gender, education, age, race etc.
\end{itemize}\medskip
Empirical Analysis for reduced form approach (there are more in the text)
\begin{itemize}\setlength\itemsep{3pt}
\item Uses two measures of displacements
\begin{itemize}\setlength\itemsep{3pt}
\item Do share of labor within an industry decline? (not just due to automation)
\item How much of the decline in labor shares are driven by automation?
\end{itemize}
\item Starts with a descriptive relationship between displacement measures and wages
\item Then see if it holds after controlling for industry and demographic characteristics
\end{itemize}\medskip
\end{frame}


\begin{frame}
\frametitle{Automation and labor shares negatively correlated} 
\begin{figure}
  \centering
  \includegraphics[keepaspectratio, width=1\textwidth]{ar(2022)_f1.png}
\end{figure}
\end{frame}

\begin{frame}
\frametitle{Middle-wage workers face highest task displacement} 
\begin{figure}
  \centering
  \includegraphics[keepaspectratio, width=0.6\textwidth]{ar(2022)_f2.png}
\end{figure}
\end{frame}

\begin{frame}
\frametitle{Task displacement negatively correlated with wage growth} 
\begin{figure}
  \centering
  \includegraphics[keepaspectratio, width=0.8\textwidth]{ar(2022)_t1.png}
\end{figure}
\end{frame}

\begin{frame}
\frametitle{Results similar with alternative displacement measure} 
\begin{figure}
  \centering
  \includegraphics[keepaspectratio, width=0.8\textwidth]{ar(2022)_t2.png}
\end{figure}
\end{frame}


\begin{frame}
\frametitle{Remaining questions on inequality} 
Reconciling differences in interpretation remains tricky
\begin{itemize}\setlength\itemsep{3pt}
\item Definition of income, tax units may differ
\item Choices on what counts leads to drastically different conclusions
\end{itemize}\medskip
How do we address inequality?
\begin{itemize}\setlength\itemsep{3pt}
\item Redistribution (transfers and taxation) could be one way
\item But it can distort incentives of both the poor and the rich \begin{itemize}\setlength\itemsep{3pt}
\item Labor incentives could be distorted for both (but for different reasons)
\end{itemize}
\item Fiscal stability concerns vs long-run gains from more equal distribution
\item What policies make opportunities more equal?
 \begin{itemize}\setlength\itemsep{3pt}
\item Early childhood care, education and health interventions, cash transfers, wealth taxes?
\end{itemize}
\end{itemize}\medskip
\end{frame}



%%%%%%%%%%%
\end{document}
